% ============================================================================
%  04-02: NORMALIZACIÓN Y REDISEÑO
% ============================================================================

\section{Normalización y redise\~no (``DESPU\'ES'')}

\subsection{Aplicaci\'on de las Formas Normales}

\subsubsection{Primera Forma Normal (1FN)}

Una relaci\'on est\'a en 1FN si todos sus atributos son at\'omicos (no hay grupos repetitivos ni atributos multivaluados).

\textbf{Violaci\'on detectada:} La columna \texttt{CarritosAbandonados.ItemsJson} almacena m\'ultiples \'items de carrito codificados como una cadena JSON en un campo \texttt{text}.
Un solo atributo contiene en realidad una lista de pares \{productoId, cantidad, precio\}, lo que constituye un grupo repetitivo encubierto.

\textbf{Correcci\'on aplicada:} Se migra la columna a tipo \texttt{jsonb} y se a\~nade un \'indice GIN para permitir b\'usquedas eficientes sobre la estructura JSON sin crear una tabla adicional.
Supabase y PostgreSQL soportan de forma nativa el tipo \texttt{jsonb} con operadores de consulta (\texttt{@>}, \texttt{?}, \texttt{jsonb\_array\_elements}).

\subsubsection{Segunda Forma Normal (2FN)}

Una relaci\'on est\'a en 2FN si est\'a en 1FN y cada atributo no clave depende funcionalmente de \textit{toda} la clave primaria (no de una parte de ella).
Esta forma normal s\'olo es relevante cuando la clave primaria es compuesta.

\textbf{Tabla analizada:} \texttt{ComboComponentes} tiene clave primaria compuesta (\texttt{ComboId}, \texttt{ComponenteId}) m\'as el atributo \texttt{Cantidad}.
El atributo \texttt{Cantidad} depende funcionalmente de ambas columnas de la clave (``cu\'antos de este componente van en este combo''), por lo que ya est\'a en 2FN.

\textbf{Tabla con atributos parcialmente dependientes:} Si \texttt{Componentes} incluyera \texttt{NombreCategoria} como texto libre junto con \texttt{CategoriaId}, \texttt{NombreCategoria} depender\'ia \'unicamente de \texttt{CategoriaId} (no de \texttt{Id} del componente), violando la 2FN en un contexto de clave compuesta hipot\'etica.
La soluci\'on es crear la tabla \texttt{Categorias} y referenciarla con \texttt{FK}.

\subsubsection{Tercera Forma Normal (3FN)}

Una relaci\'on est\'a en 3FN si est\'a en 2FN y ning\'un atributo no clave depende transitivamente de la clave primaria (es decir, no hay dependencias entre atributos no clave).

\textbf{Violaci\'on en \texttt{AspNetUsers}:}
\begin{itemize}
  \item \texttt{Id} $\rightarrow$ \texttt{NombreCompleto} (dependencia directa, correcta).
  \item \texttt{Id} $\rightarrow$ \texttt{FullName} (dependencia transitiva: \texttt{NombreCompleto} $\rightarrow$ \texttt{FullName}).
\end{itemize}
La columna \texttt{FullName} es una copia de \texttt{NombreCompleto}; eliminarla resuelve la dependencia transitiva.

\textbf{Violaci\'on en \texttt{Componentes}:}
\begin{itemize}
  \item \texttt{Id} $\rightarrow$ \texttt{Tipo} (texto libre, e.g. ``Tarjetas de Video'').
  \item \texttt{Tipo} $\rightarrow$ \texttt{NombreCategoria} (si se a\~nadiera esa columna, ser\'ia transitiva).
\end{itemize}
La soluci\'on es separar \texttt{Tipo} en una tabla \texttt{Categorias} (\texttt{Id}, \texttt{Nombre}) y reemplazar \texttt{Tipo text} por \texttt{CategoriaId integer FK}.
An\'alogamente, \texttt{Marca text} se reemplaza por \texttt{MarcaId integer FK} referenciando la tabla \texttt{Marcas}.

\subsection{Dise\~no del esquema DESPU\'ES}

\subsubsection{Separaci\'on en esquemas de negocio}

Se organiza la base de datos en cinco esquemas especializados:

\begin{center}
\begin{tablaAPA}
\begin{tabular}{L{2.8cm} L{5.5cm} L{5.5cm}}
\toprule
\textbf{Esquema} & \textbf{Tablas} & \textbf{Prop\'osito} \\
\midrule
\texttt{catalogo} & Componentes, Categorias, Marcas, Combos, ComboComponentes & Cat\'alogo de productos y agrupaciones \\
\texttt{ventas} & Pedidos, PedidoItems, CarritosAbandonados & Ciclo de vida de ventas \\
\texttt{operaciones} & Sucursales, Proveedores, Regiones & Log\'istica y abastecimiento \\
\texttt{seguridad} & AspNetUsers (extendida) + tablas Identity & Identidad y autorizaci\'on \\
\texttt{auditoria} & AuditoriaPedidos, AuditoriaComponentes & Trazabilidad de cambios \\
\bottomrule
\end{tabular}
\end{tablaAPA}
\captionof{table}{Distribuci\'on de tablas por esquema --- Estado DESPU\'ES}
\label{tab:esquemas}
\end{center}

\subsection{Diccionario de datos completo}

\subsubsection{Esquema \texttt{catalogo} --- Tabla \texttt{Categorias}}

\begin{center}
\begin{tablaAPA}
\begin{tabular}{L{2.8cm} L{2.0cm} C{1.4cm} L{7.0cm}}
\toprule
\textbf{Columna} & \textbf{Tipo} & \textbf{Nulo} & \textbf{Restricciones / Descripci\'on} \\
\midrule
Id      & serial      & NO & PK, autoincremental \\
Nombre  & varchar(60) & NO & UNIQUE; nombre de la categor\'ia (ej. ``Tarjetas de Video'') \\
Slug    & varchar(80) & NO & UNIQUE; identificador URL-safe (ej. ``tarjetas-de-video'') \\
\bottomrule
\end{tabular}
\end{tablaAPA}
\captionof{table}{Diccionario de datos --- \texttt{catalogo.Categorias}}
\label{tab:dict-categorias}
\end{center}

\subsubsection{Esquema \texttt{catalogo} --- Tabla \texttt{Marcas}}

\begin{center}
\begin{tablaAPA}
\begin{tabular}{L{2.8cm} L{2.0cm} C{1.4cm} L{7.0cm}}
\toprule
\textbf{Columna} & \textbf{Tipo} & \textbf{Nulo} & \textbf{Restricciones / Descripci\'on} \\
\midrule
Id      & serial      & NO & PK, autoincremental \\
Nombre  & varchar(80) & NO & UNIQUE; nombre de la marca (ej. ``NVIDIA'', ``ASUS ROG'') \\
Pa\'is   & varchar(60) & S\'I & Pa\'is de origen de la marca \\
\bottomrule
\end{tabular}
\end{tablaAPA}
\captionof{table}{Diccionario de datos --- \texttt{catalogo.Marcas}}
\label{tab:dict-marcas}
\end{center}

\subsubsection{Esquema \texttt{catalogo} --- Tabla \texttt{Componentes}}

\begin{center}
\begin{tablaAPA}
\begin{tabular}{L{2.8cm} L{2.2cm} C{1.2cm} L{7.0cm}}
\toprule
\textbf{Columna} & \textbf{Tipo} & \textbf{Nulo} & \textbf{Restricciones / Descripci\'on} \\
\midrule
Id          & serial         & NO & PK \\
Nombre      & varchar(200)   & NO & Nombre del componente \\
Descripcion & text           & S\'I & Descripci\'on t\'ecnica \\
CategoriaId & integer        & NO & FK $\to$ \texttt{Categorias.Id} ON DELETE RESTRICT \\
MarcaId     & integer        & NO & FK $\to$ \texttt{Marcas.Id} ON DELETE RESTRICT \\
Precio      & numeric(10,2)  & NO & CHECK (Precio $>$ 0) \\
Stock       & integer        & NO & CHECK (Stock $\geq$ 0); DEFAULT 0 \\
ImagenUrl   & varchar(500)   & S\'I & URL de imagen del producto \\
Activo      & boolean        & NO & DEFAULT true; habilita/deshabilita el producto \\
CreadoEn    & timestamptz    & NO & DEFAULT now() \\
\bottomrule
\end{tabular}
\end{tablaAPA}
\captionof{table}{Diccionario de datos --- \texttt{catalogo.Componentes} (Estado DESPU\'ES)}
\label{tab:dict-componentes}
\end{center}

\subsubsection{Esquema \texttt{ventas} --- Tabla \texttt{PedidoItems} (redise\~nada)}

\begin{center}
\begin{tablaAPA}
\begin{tabular}{L{2.8cm} L{2.2cm} C{1.2cm} L{7.0cm}}
\toprule
\textbf{Columna} & \textbf{Tipo} & \textbf{Nulo} & \textbf{Restricciones / Descripci\'on} \\
\midrule
Id            & serial        & NO & PK \\
PedidoId      & integer       & NO & FK $\to$ \texttt{Pedidos.Id} ON DELETE CASCADE \\
ComponenteId  & integer       & S\'I & FK $\to$ \texttt{Componentes.Id} ON DELETE RESTRICT; nulo si es Combo \\
ComboId       & integer       & S\'I & FK $\to$ \texttt{Combos.Id} ON DELETE RESTRICT; nulo si es Componente \\
Cantidad      & integer       & NO & CHECK (Cantidad $>$ 0) \\
PrecioUnitario& numeric(10,2) & NO & CHECK (PrecioUnitario $>$ 0); precio al momento de la compra \\
\bottomrule
\end{tabular}
\end{tablaAPA}
\captionof{table}{Diccionario de datos --- \texttt{ventas.PedidoItems} (Estado DESPU\'ES)}
\label{tab:dict-pedidoitems}
\end{center}

La restricci\'on \texttt{CHECK ((ComponenteId IS NOT NULL AND ComboId IS NULL) OR (ComboId IS NOT NULL AND ComponenteId IS NULL))} garantiza que cada \'item referencie exactamente un tipo de producto, resolviendo la referencia polim\'orfica sin FK del estado anterior.

\subsection{Diagrama Entidad-Relaci\'on --- Estado DESPU\'ES}

\figuraAPA{fig:der-despues}{Diagrama Entidad-Relaci\'on normalizado --- Estado DESPU\'ES}{%
\begin{tikzpicture}[
  ent/.style={draw, rectangle, rounded corners=3pt, fill=green!12, minimum width=2.6cm, minimum height=0.65cm, font=\small\bfseries, align=center},
  aux/.style={draw, rectangle, rounded corners=3pt, fill=blue!10, minimum width=2.2cm, minimum height=0.6cm, font=\small, align=center},
  aud/.style={draw, rectangle, rounded corners=3pt, fill=gray!15, minimum width=2.4cm, minimum height=0.6cm, font=\small, align=center},
  fk/.style={-Stealth, thick},
  node distance=1.3cm and 2.0cm
]
% Esquema catalogo
\node[aux] (cat) {Categorias};
\node[aux, right=2cm of cat] (mar) {Marcas};
\node[ent, below=1.2cm of cat] (comp) {Componentes};
\node[ent, right=2.3cm of comp] (combos) {Combos};
\node[ent, right=2.3cm of combos] (ccomp) {Combo\\Componentes};

% Esquema ventas
\node[ent, below=1.8cm of comp] (ped) {Pedidos};
\node[ent, right=2.3cm of ped] (pitems) {PedidoItems};
\node[ent, right=2.3cm of pitems] (carr) {Carritos\\Abandonados};

% Esquema operaciones
\node[aux, below=1.6cm of ped] (suc) {Sucursales};
\node[aux, right=2.2cm of suc] (prov) {Proveedores};
\node[aux, left=2.2cm of suc] (reg) {Regiones};

% Esquema seguridad
\node[aux, left=2.5cm of ped] (users) {AspNetUsers};

% Esquema auditoria
\node[aud, below=1.6cm of pitems] (audped) {Auditoria\\Pedidos};
\node[aud, right=1.5cm of audped] (audcomp) {Auditoria\\Componentes};

% FKs
\draw[fk] (comp) -- node[left,font=\tiny]{FK} (cat);
\draw[fk] (comp.north east) -- node[above,font=\tiny]{FK} (mar.south west);
\draw[fk] (ccomp) -- node[above,font=\tiny]{FK} (combos);
\draw[fk] (ccomp.south) -- node[right,font=\tiny]{FK} (comp.north east);
\draw[fk] (ped) -- node[above,font=\tiny]{FK} (users);
\draw[fk] (pitems) -- node[above,font=\tiny]{FK} (ped);
\draw[fk] (pitems.north) -- node[right,font=\tiny]{FK opt} (comp.south east);
\draw[fk] (pitems.north east) -- node[above,font=\tiny]{FK opt} (combos.south);
\draw[fk] (carr) -- node[right,font=\tiny]{FK} (users);
\draw[fk] (users) -- node[left,font=\tiny]{FK} (suc);
\draw[fk] (users.south west) -- node[below,font=\tiny]{FK} (reg.north east);
\draw[fk] (audped) -- node[left,font=\tiny]{trig} (ped);
\draw[fk] (audcomp) -- node[right,font=\tiny]{trig} (comp.south);
\end{tikzpicture}%
}{Elaboraci\'on propia basada en el redise\~no propuesto para la materia Base de Datos II}
